\documentclass{beamer}
\usetheme{Madrid}
\usecolortheme{beaver} % A nice dark red theme fitting for steel/metallurgy

\title{The Physics of Steel}
\subtitle{Governing Differential Equations: From Furnace to Failure}
\author{Metallurgy \& Physics}
\date{\today}

\begin{document}

\begin{frame}
    \titlepage
\end{frame}

\begin{frame}{Outline}
    \tableofcontents
\end{frame}

% --- SECTION 1: FORMATION ---
\section{1. Formation: Heat and Phase Transformations}

\begin{frame}{Formation: Heat Transfer}
    \textbf{The Heat Equation} \\
    Controlling the cooling rate is critical to determining the final crystalline structure (e.g., Martensite vs. Pearlite). The temperature distribution $T$ over time $t$ is governed by:

    $$\rho C_p \frac{\partial T}{\partial t} = \nabla \cdot (k \nabla T) + \dot{Q}$$

    \begin{itemize}
        \item $\rho$: Density
        \item $C_p$: Specific heat capacity
        \item $k$: Thermal conductivity
        \item $\dot{Q}$: Internal heat generation
    \end{itemize}
\end{frame}

\begin{frame}{Formation: Carbon Diffusion}
    \textbf{Fick's Second Law} \\
    The migration of Carbon atoms through the Iron lattice in real space:
    $$\frac{\partial C}{\partial t} = D \frac{\partial^2 C}{\partial x^2}$$

    \vspace{0.5cm}
    \textbf{Fourier Space Representation} \\
    Transforming to the frequency domain $k$ reveals that diffusion acts as a low-pass filter, rapidly smoothing out sharp concentration gradients:
    $$\frac{\partial \tilde{C}(k,t)}{\partial t} = -D k^2 \tilde{C}(k,t)$$
\end{frame}

\begin{frame}{Formation: Diffusion in a Periodic Lattice}
    \textbf{The Smoluchowski Equation} \\
    When accounting for the periodic potential $V_{Fe}$ of the crystalline grid, diffusion includes a drift term:
    $$\frac{\partial C}{\partial t} = \nabla \cdot \left( D \nabla C + \frac{D}{k_B T} C (\nabla V_{Fe}) \right)$$

    \vspace{0.5cm}
    \textbf{The Low-Temperature Limit} \\
    As $T \to 0$, the equilibrium distribution of Carbon atoms becomes strictly locked into the interstitial minima, mathematically forming a periodic Dirac comb:
    $$C(x) \propto \lim_{T \to 0} \exp\left[ \frac{V_{amp}}{k_B T} \cos\left(\frac{2\pi x}{L}\right) \right]$$
\end{frame}

\begin{frame}{Formation: Phase Transformation Kinetics}
    \textbf{The JMAK Equation} \\
    The Johnson-Mehl-Avrami-Kolmogorov (JMAK) equation tracks the volumetric fraction $f$ of the steel that has transformed into a new phase:

    $$\frac{df}{dt} = n \cdot K \cdot t^{n-1} (1 - f)$$

    \begin{itemize}
        \item $f$: Fraction of transformed material
        \item $n$: Avrami exponent (depends on nucleation type)
        \item $K$: Temperature-dependent rate constant
    \end{itemize}
\end{frame}

% --- SECTION 2: MICROSTRUCTURE ---
\section{2. Microstructure Evolution}

\begin{frame}{Microstructure: Phase-Field Models}
    \textbf{Allen-Cahn / Cahn-Hilliard Equations} \\
    To model the physical shape and evolution of grains (like dendrites) or spinodal decomposition, we track a continuous phase-field variable $\phi$. The system evolves to minimize total free energy $F$:

    $$\frac{\partial \phi}{\partial t} = -M \left( \frac{\delta F}{\delta \phi} \right)$$

    \begin{itemize}
        \item $\phi$: Phase-field order parameter
        \item $M$: Mobility parameter
        \item $\frac{\delta F}{\delta \phi}$: Variational derivative of the free energy functional
    \end{itemize}
\end{frame}

% --- SECTION 3: DEGRADATION ---
\section{3. Degradation: Corrosion and Fatigue}

\begin{frame}{Degradation: Electrochemical Corrosion}
    \textbf{The Butler-Volmer Equation} \\
    Corrosion is an electrochemical decay process driven by overpotential $\eta$. The rate of metal loss (current density $j$) is given by:

    $$j = j_0 \left\{ \exp \left[ \frac{\alpha_a z F \eta}{RT} \right] - \exp \left[ -\frac{\alpha_c z F \eta}{RT} \right] \right\}$$

    \begin{itemize}
        \item $j_0$: Exchange current density
        \item $\alpha_a, \alpha_c$: Anodic and cathodic charge transfer coefficients
        \item $\eta$: Overpotential (driving force)
    \end{itemize}
\end{frame}

\begin{frame}{Degradation: Mechanical Fatigue}
    \textbf{Paris' Law} \\
    Under cyclic mechanical loading, microscopic defects grow into macroscopic fractures. The crack growth rate per load cycle $N$ is governed by an ordinary differential equation:

    $$\frac{da}{dN} = C(\Delta K)^m$$

    \begin{itemize}
        \item $a$: Crack length
        \item $N$: Number of load cycles
        \item $\Delta K$: Stress intensity factor range
        \item $C, m$: Material constants
    \end{itemize}
\end{frame}

\begin{frame}{Summary}
    \begin{center}
        \Large Steel is a dynamic material governed by a continuous balance of thermodynamics and kinetics. \\
        \vspace{1cm}
        From the high-temperature PDEs of diffusion and heat flow to the cyclic ODEs of structural fatigue, differential equations define its entire life cycle.
    \end{center}
\end{frame}

\end{document}